\subsection{Diagonal equality}

\begin{proposition}[Diagonal fiber]\label{thm:diagFiber}\lean{Obsmetrologia.diagFiber}\leanok\uses{thm:predication, thm:existsT, thm:existsF}
The diagonal fiber over $(s,o)$ is a singleton when $s=o$ and empty otherwise.
It is the finite fiber along which Lawvere's diagonal-$\Sigma$ construction
builds equality.
\end{proposition}

\begin{proposition}[Classical and paraconsistent equality]\label{thm:classicalEq}\lean{Obsmetrologia.classicalEq, Obsmetrologia.paraconsistentEq}\leanok\uses{thm:diagFiber, thm:existsT, thm:existsF, thm:charge-joinTList, thm:charge-joinList}
Classical equality is the truth-order existential image of the diagonal
$t$-terminal. Paraconsistent equality is the knowledge-order existential image
of the diagonal $\top$-terminal. Equality is therefore constructed in each
order, not postulated.
\end{proposition}

\begin{theorem}[Classical equality is two-valued]\label{thm:classicalEq-eq}\lean{Obsmetrologia.classicalEq_eq}\leanok\uses{thm:classicalEq, thm:existsT-adjunction}
Classical equality returns $t$ on the diagonal and $f$ off it.
\end{theorem}

\begin{theorem}[Paraconsistent equality is $\top$/$\bot$]\label{thm:paraconsistentEq-eq}\lean{Obsmetrologia.paraconsistentEq_eq}\leanok\uses{thm:classicalEq, thm:existsT-adjunction, thm:four-bilattice}
Paraconsistent equality returns $\top$ on the diagonal and $\bot$ off it.
\end{theorem}

\begin{corollary}[Reflexivity in both orders]\label{cor:eq-refl}\lean{Obsmetrologia.classicalEq_refl, Obsmetrologia.paraconsistentEq_refl}\leanok\uses{thm:classicalEq-eq, thm:paraconsistentEq-eq}
Both equalities are reflexive, but they carry different charges on the same
diagonal: $t$ for classical truth-order equality and $\top$ for
knowledge-order paraconsistent equality.
\end{corollary}
